% --- Archon physics formalization source begin ---
% archon:physics
% archon:covers PhyXMiniProblems/problem_phyx_mini_0866.lean
% archon:source-report reports/phyx_mini/problem_phyx_mini_0866.source.json

\chapter{Physics problem phyx\_mini\_0866}
\label{ch:PhyXMiniProblems_problem_phyx_mini_0866}

\paragraph{Problem source.}
\#\# Physical scenario

In the figure, a proton is fired with a speed of 200,000 m/s from the midpoint of the capacitor toward the positive plate.

\#\# Question

What is the proton’s speed as it collides with the negative plate?

\#\# Answer choices

- A: A: 1.59\textbackslash{}times 10\textasciicircum{}\{5\}m/s

- B: B: 1.49\textbackslash{}times 10\textasciicircum{}\{5\}m/s

- C: C: 2.19\textbackslash{}times 10\textasciicircum{}\{5\}m/s

- D: D: 2.96\textbackslash{}times 10\textasciicircum{}\{5\}m/s

\#\# Figure caption (auxiliary; use the image as primary evidence)

The image depicts a diagram illustrating the motion of a positively charged particle between two parallel plates. 

- **Plates**: 
  - The left plate is labeled "0 V" and has four small horizontal lines.
  - The right plate is labeled "500 V" and contains four plus signs, indicating a positive charge.

- **Particle**: 
  - There is a red circle with a plus sign in the center, representing a positive charge.
  - The particle is moving from the right plate to the left plate.

- **Arrow**:
  - A blue curved arrow indicates the path and direction of the particle's motion from the positive plate to the neutral plate.

- **Text**: 
  - "200,000 m/s" is written next to the particle, indicating its speed.

Overall, this figure represents the movement of a positively charged particle in an electric field created by the voltage difference between the two plates.

\#\# Dataset metadata

Electrostatics; Physical Model Grounding Reasoning, Predictive Reasoning

\paragraph{Recorded answer/context.}
D: D: 2.96\textbackslash{}times 10\textasciicircum{}\{5\}m/s

\paragraph{Figure/image path.}
/root/proposal\_for\_physic/science-mango/phyx\_mini\_run/phyx\_data/test\_image/866.png

\paragraph{Formalization target.}
create a compiling Lean file with sorry bodies at `PhyXMiniProblems/problem\_phyx\_mini\_0866.lean`. The Lean declarations must preserve the physical quantities, dimensions or dimensional roles, figure labels, governing-law hypotheses, and final relation expressed by this problem.
Use Mathlib/Physlib names found through LeanExplore where available. If a physics API is missing, introduce faithful local abstractions rather than scalar placeholder aliases.

\begin{theorem}[Physics formalization target]
\label{thm:physics:phyx_mini_0866:target}
\lean{PhyXMiniProblems.ProblemPhyXMini0866.proton_collision_speed_matches_choice_D}
\uses{def:physics:phyx-mini-0866:phyxminiproblems-problemphyxmini0866-massinkilograms, def:physics:phyx-mini-0866:phyxminiproblems-problemphyxmini0866-chargemagnitudeincoulombs, def:physics:phyx-mini-0866:phyxminiproblems-problemphyxmini0866-speedinmeterspersecond, def:physics:phyx-mini-0866:phyxminiproblems-problemphyxmini0866-electricpotentialinvolts, def:physics:phyx-mini-0866:phyxminiproblems-problemphyxmini0866-capacitorprotonsetup, def:physics:phyx-mini-0866:phyxminiproblems-problemphyxmini0866-matchescapacitorprotonscenario, def:physics:phyx-mini-0866:phyxminiproblems-problemphyxmini0866-matchessuppliedcapacitorfigureanddata, def:physics:phyx-mini-0866:phyxminiproblems-problemphyxmini0866-hasstandardprotoncalibration, def:physics:phyx-mini-0866:phyxminiproblems-problemphyxmini0866-hasphysicalcapacitorprotonparameters, def:physics:phyx-mini-0866:phyxminiproblems-problemphyxmini0866-satisfiesuniformparallelplatepotentiallaw, def:physics:phyx-mini-0866:phyxminiproblems-problemphyxmini0866-satisfieselectrostaticenergylaws, def:physics:phyx-mini-0866:phyxminiproblems-problemphyxmini0866-displayedspeedinmeterspersecond, def:physics:phyx-mini-0866:phyxminiproblems-problemphyxmini0866-roundstodisplayedspeed, def:physics:phyx-mini-0866:phyxminiproblems-problemphyxmini0866-isclosestdisplayedspeed}
The assigned autoformalize agent should translate this physics problem into Lean declarations in the covered file, with theorem and lemma proof bodies written as `by sorry`.
\end{theorem}
\begin{proof}
This is an autoformalization task, not a proof task. Produce statements that can later be proved by the physics prover without weakening the physical model.
\end{proof}

% --- Archon declaration topology begin ---
\section{Declaration topology}
\begin{definition}[Mass Quantity]
  \label{def:physics:phyx-mini-0866:phyxminiproblems-problemphyxmini0866-massquantity}
  \lean{PhyXMiniProblems.ProblemPhyXMini0866.MassQuantity}
  A nonnegative physical mass
\end{definition}
\begin{proof}
  This is the declared model definition of Mass Quantity.
\end{proof}
\begin{definition}[Charge Magnitude Quantity]
  \label{def:physics:phyx-mini-0866:phyxminiproblems-problemphyxmini0866-chargemagnitudequantity}
  \lean{PhyXMiniProblems.ProblemPhyXMini0866.ChargeMagnitudeQuantity}
  A nonnegative electric-charge magnitude
\end{definition}
\begin{proof}
  This is the declared model definition of Charge Magnitude Quantity.
\end{proof}
\begin{definition}[speed Dimension]
  \label{def:physics:phyx-mini-0866:phyxminiproblems-problemphyxmini0866-speeddimension}
  \lean{PhyXMiniProblems.ProblemPhyXMini0866.speedDimension}
  The physical dimension L T⁻¹ of speed
\end{definition}
\begin{proof}
  This is the declared model definition of speed Dimension.
\end{proof}
\begin{definition}[Speed Quantity]
  \label{def:physics:phyx-mini-0866:phyxminiproblems-problemphyxmini0866-speedquantity}
  \lean{PhyXMiniProblems.ProblemPhyXMini0866.SpeedQuantity}
  \uses{def:physics:phyx-mini-0866:phyxminiproblems-problemphyxmini0866-speeddimension}
  A nonnegative physical speed magnitude
\end{definition}
\begin{proof}
  This is the declared model definition of Speed Quantity.
\end{proof}
\begin{definition}[electric Potential Dimension]
  \label{def:physics:phyx-mini-0866:phyxminiproblems-problemphyxmini0866-electricpotentialdimension}
  \lean{PhyXMiniProblems.ProblemPhyXMini0866.electricPotentialDimension}
  The dimension M L² T⁻² C⁻¹ of electric potential
\end{definition}
\begin{proof}
  This is the declared model definition of electric Potential Dimension.
\end{proof}
\begin{definition}[Electric Potential Quantity]
  \label{def:physics:phyx-mini-0866:phyxminiproblems-problemphyxmini0866-electricpotentialquantity}
  \lean{PhyXMiniProblems.ProblemPhyXMini0866.ElectricPotentialQuantity}
  \uses{def:physics:phyx-mini-0866:phyxminiproblems-problemphyxmini0866-electricpotentialdimension}
  A signed, unit-independent physical electric potential
\end{definition}
\begin{proof}
  This is the declared model definition of Electric Potential Quantity.
\end{proof}
\begin{definition}[mass In Kilograms]
  \label{def:physics:phyx-mini-0866:phyxminiproblems-problemphyxmini0866-massinkilograms}
  \lean{PhyXMiniProblems.ProblemPhyXMini0866.massInKilograms}
  \uses{def:physics:phyx-mini-0866:phyxminiproblems-problemphyxmini0866-massquantity}
  SI mass readout in kilograms
\end{definition}
\begin{proof}
  This is the declared model definition of mass In Kilograms.
\end{proof}
\begin{definition}[charge Magnitude In Coulombs]
  \label{def:physics:phyx-mini-0866:phyxminiproblems-problemphyxmini0866-chargemagnitudeincoulombs}
  \lean{PhyXMiniProblems.ProblemPhyXMini0866.chargeMagnitudeInCoulombs}
  \uses{def:physics:phyx-mini-0866:phyxminiproblems-problemphyxmini0866-chargemagnitudequantity}
  SI charge-magnitude readout in coulombs
\end{definition}
\begin{proof}
  This is the declared model definition of charge Magnitude In Coulombs.
\end{proof}
\begin{definition}[speed In Meters Per Second]
  \label{def:physics:phyx-mini-0866:phyxminiproblems-problemphyxmini0866-speedinmeterspersecond}
  \lean{PhyXMiniProblems.ProblemPhyXMini0866.speedInMetersPerSecond}
  \uses{def:physics:phyx-mini-0866:phyxminiproblems-problemphyxmini0866-speedquantity}
  SI speed readout in metres per second
\end{definition}
\begin{proof}
  This is the declared model definition of speed In Meters Per Second.
\end{proof}
\begin{definition}[electric Potential In Volts]
  \label{def:physics:phyx-mini-0866:phyxminiproblems-problemphyxmini0866-electricpotentialinvolts}
  \lean{PhyXMiniProblems.ProblemPhyXMini0866.electricPotentialInVolts}
  \uses{def:physics:phyx-mini-0866:phyxminiproblems-problemphyxmini0866-electricpotentialquantity}
  SI electric-potential readout in volts
\end{definition}
\begin{proof}
  This is the declared model definition of electric Potential In Volts.
\end{proof}
\begin{definition}[energy In Joules]
  \label{def:physics:phyx-mini-0866:phyxminiproblems-problemphyxmini0866-energyinjoules}
  \lean{PhyXMiniProblems.ProblemPhyXMini0866.energyInJoules}
  SI energy readout in joules
\end{definition}
\begin{proof}
  This is the declared model definition of energy In Joules.
\end{proof}
\begin{definition}[elementary Charge In Coulombs]
  \label{def:physics:phyx-mini-0866:phyxminiproblems-problemphyxmini0866-elementarychargeincoulombs}
  \lean{PhyXMiniProblems.ProblemPhyXMini0866.elementaryChargeInCoulombs}
  The exact elementary-charge magnitude in coulombs, grounded in Physlib's ChargeUnit.elementaryCharge rather than introduced as an untyped decimal
\end{definition}
\begin{proof}
  This is the declared model definition of elementary Charge In Coulombs.
\end{proof}
\begin{definition}[standard Proton Mass In Kilograms]
  \label{def:physics:phyx-mini-0866:phyxminiproblems-problemphyxmini0866-standardprotonmassinkilograms}
  \lean{PhyXMiniProblems.ProblemPhyXMini0866.standardProtonMassInKilograms}
  CODATA proton mass readout used by the problem's nonrelativistic model
\end{definition}
\begin{proof}
  This is the declared model definition of standard Proton Mass In Kilograms.
\end{proof}
\begin{definition}[Plate Label]
  \label{def:physics:phyx-mini-0866:phyxminiproblems-problemphyxmini0866-platelabel}
  \lean{PhyXMiniProblems.ProblemPhyXMini0866.PlateLabel}
  The two capacitor plates as placed in the primary bitmap
\end{definition}
\begin{proof}
  This is the declared model definition of Plate Label.
\end{proof}
\begin{definition}[Plate Polarity]
  \label{def:physics:phyx-mini-0866:phyxminiproblems-problemphyxmini0866-platepolarity}
  \lean{PhyXMiniProblems.ProblemPhyXMini0866.PlatePolarity}
  The electrostatic role of a capacitor plate
\end{definition}
\begin{proof}
  This is the declared model definition of Plate Polarity.
\end{proof}
\begin{definition}[Particle Species]
  \label{def:physics:phyx-mini-0866:phyxminiproblems-problemphyxmini0866-particlespecies}
  \lean{PhyXMiniProblems.ProblemPhyXMini0866.ParticleSpecies}
  Particle species distinguished by this model
\end{definition}
\begin{proof}
  This is the declared model definition of Particle Species.
\end{proof}
\begin{definition}[Charge Sign]
  \label{def:physics:phyx-mini-0866:phyxminiproblems-problemphyxmini0866-chargesign}
  \lean{PhyXMiniProblems.ProblemPhyXMini0866.ChargeSign}
  Sign of a particle's electric charge
\end{definition}
\begin{proof}
  This is the declared model definition of Charge Sign.
\end{proof}
\begin{definition}[Plate Arrangement]
  \label{def:physics:phyx-mini-0866:phyxminiproblems-problemphyxmini0866-platearrangement}
  \lean{PhyXMiniProblems.ProblemPhyXMini0866.PlateArrangement}
  Arrangement of the two long capacitor plates
\end{definition}
\begin{proof}
  This is the declared model definition of Plate Arrangement.
\end{proof}
\begin{definition}[Horizontal Direction]
  \label{def:physics:phyx-mini-0866:phyxminiproblems-problemphyxmini0866-horizontaldirection}
  \lean{PhyXMiniProblems.ProblemPhyXMini0866.HorizontalDirection}
  Horizontal directions named by the side-view diagram
\end{definition}
\begin{proof}
  This is the declared model definition of Horizontal Direction.
\end{proof}
\begin{definition}[Motion Event]
  \label{def:physics:phyx-mini-0866:phyxminiproblems-problemphyxmini0866-motionevent}
  \lean{PhyXMiniProblems.ProblemPhyXMini0866.MotionEvent}
  Named instants in the proton's one-dimensional motion
\end{definition}
\begin{proof}
  This is the declared model definition of Motion Event.
\end{proof}
\begin{definition}[Across Gap Location]
  \label{def:physics:phyx-mini-0866:phyxminiproblems-problemphyxmini0866-acrossgaplocation}
  \lean{PhyXMiniProblems.ProblemPhyXMini0866.AcrossGapLocation}
  Qualitative locations across the gap between the plates
\end{definition}
\begin{proof}
  This is the declared model definition of Across Gap Location.
\end{proof}
\begin{definition}[Drawn Plate Mark]
  \label{def:physics:phyx-mini-0866:phyxminiproblems-problemphyxmini0866-drawnplatemark}
  \lean{PhyXMiniProblems.ProblemPhyXMini0866.DrawnPlateMark}
  The charge marks drawn repeatedly inside a plate
\end{definition}
\begin{proof}
  This is the declared model definition of Drawn Plate Mark.
\end{proof}
\begin{definition}[Capacitor Proton Figure]
  \label{def:physics:phyx-mini-0866:phyxminiproblems-problemphyxmini0866-capacitorprotonfigure}
  \lean{PhyXMiniProblems.ProblemPhyXMini0866.CapacitorProtonFigure}
  \uses{def:physics:phyx-mini-0866:phyxminiproblems-problemphyxmini0866-platelabel, def:physics:phyx-mini-0866:phyxminiproblems-problemphyxmini0866-horizontaldirection, def:physics:phyx-mini-0866:phyxminiproblems-problemphyxmini0866-drawnplatemark}
  Literal features of phyx\_data/test\_image/866.png. The real-valued horizontal coordinates are diagram coordinates only; they are not physical lengths
\end{definition}
\begin{proof}
  This is the declared model definition of Capacitor Proton Figure.
\end{proof}
\begin{definition}[Capacitor Proton Setup]
  \label{def:physics:phyx-mini-0866:phyxminiproblems-problemphyxmini0866-capacitorprotonsetup}
  \lean{PhyXMiniProblems.ProblemPhyXMini0866.CapacitorProtonSetup}
  \uses{def:physics:phyx-mini-0866:phyxminiproblems-problemphyxmini0866-massquantity, def:physics:phyx-mini-0866:phyxminiproblems-problemphyxmini0866-chargemagnitudequantity, def:physics:phyx-mini-0866:phyxminiproblems-problemphyxmini0866-speedquantity, def:physics:phyx-mini-0866:phyxminiproblems-problemphyxmini0866-electricpotentialquantity, def:physics:phyx-mini-0866:phyxminiproblems-problemphyxmini0866-platelabel, def:physics:phyx-mini-0866:phyxminiproblems-problemphyxmini0866-platepolarity, def:physics:phyx-mini-0866:phyxminiproblems-problemphyxmini0866-particlespecies, def:physics:phyx-mini-0866:phyxminiproblems-problemphyxmini0866-chargesign, def:physics:phyx-mini-0866:phyxminiproblems-problemphyxmini0866-platearrangement, def:physics:phyx-mini-0866:phyxminiproblems-problemphyxmini0866-horizontaldirection, def:physics:phyx-mini-0866:phyxminiproblems-problemphyxmini0866-motionevent, def:physics:phyx-mini-0866:phyxminiproblems-problemphyxmini0866-acrossgaplocation, def:physics:phyx-mini-0866:phyxminiproblems-problemphyxmini0866-capacitorprotonfigure}
  Independent physical observables for the capacitor experiment. Neither the collision speed nor any displayed answer is fixed by this structure
\end{definition}
\begin{proof}
  This is the declared model definition of Capacitor Proton Setup.
\end{proof}
\begin{definition}[Matches Capacitor Proton Scenario]
  \label{def:physics:phyx-mini-0866:phyxminiproblems-problemphyxmini0866-matchescapacitorprotonscenario}
  \lean{PhyXMiniProblems.ProblemPhyXMini0866.MatchesCapacitorProtonScenario}
  \uses{def:physics:phyx-mini-0866:phyxminiproblems-problemphyxmini0866-speedinmeterspersecond, def:physics:phyx-mini-0866:phyxminiproblems-problemphyxmini0866-capacitorprotonsetup}
  Qualitative physical roles stated in the prose and represented by the curved trajectory. The zero speed at the turning point expresses reversal of this one-dimensional motion; it does not constrain the requested collision speed
\end{definition}
\begin{proof}
  This is the declared model definition of Matches Capacitor Proton Scenario.
\end{proof}
\begin{definition}[Matches Supplied Capacitor Figure And Data]
  \label{def:physics:phyx-mini-0866:phyxminiproblems-problemphyxmini0866-matchessuppliedcapacitorfigureanddata}
  \lean{PhyXMiniProblems.ProblemPhyXMini0866.MatchesSuppliedCapacitorFigureAndData}
  \uses{def:physics:phyx-mini-0866:phyxminiproblems-problemphyxmini0866-speedinmeterspersecond, def:physics:phyx-mini-0866:phyxminiproblems-problemphyxmini0866-electricpotentialinvolts, def:physics:phyx-mini-0866:phyxminiproblems-problemphyxmini0866-capacitorprotonsetup}
  Readouts and visual facts transcribed from the primary bitmap. The bitmap has five cyan minus marks on the left plate and five red plus marks on the right. No collision-speed number occurs here
\end{definition}
\begin{proof}
  This is the declared model definition of Matches Supplied Capacitor Figure And Data.
\end{proof}
\begin{definition}[Has Standard Proton Calibration]
  \label{def:physics:phyx-mini-0866:phyxminiproblems-problemphyxmini0866-hasstandardprotoncalibration}
  \lean{PhyXMiniProblems.ProblemPhyXMini0866.HasStandardProtonCalibration}
  \uses{def:physics:phyx-mini-0866:phyxminiproblems-problemphyxmini0866-massinkilograms, def:physics:phyx-mini-0866:phyxminiproblems-problemphyxmini0866-chargemagnitudeincoulombs, def:physics:phyx-mini-0866:phyxminiproblems-problemphyxmini0866-elementarychargeincoulombs, def:physics:phyx-mini-0866:phyxminiproblems-problemphyxmini0866-standardprotonmassinkilograms, def:physics:phyx-mini-0866:phyxminiproblems-problemphyxmini0866-capacitorprotonsetup}
  Standard physical calibration of the proton mass and positive elementary charge. These data do not mention the unknown collision speed
\end{definition}
\begin{proof}
  This is the declared model definition of Has Standard Proton Calibration.
\end{proof}
\begin{definition}[Has Physical Capacitor Proton Parameters]
  \label{def:physics:phyx-mini-0866:phyxminiproblems-problemphyxmini0866-hasphysicalcapacitorprotonparameters}
  \lean{PhyXMiniProblems.ProblemPhyXMini0866.HasPhysicalCapacitorProtonParameters}
  \uses{def:physics:phyx-mini-0866:phyxminiproblems-problemphyxmini0866-massinkilograms, def:physics:phyx-mini-0866:phyxminiproblems-problemphyxmini0866-chargemagnitudeincoulombs, def:physics:phyx-mini-0866:phyxminiproblems-problemphyxmini0866-speedinmeterspersecond, def:physics:phyx-mini-0866:phyxminiproblems-problemphyxmini0866-capacitorprotonsetup}
  Positivity and nondegeneracy conditions for the physical experiment
\end{definition}
\begin{proof}
  This is the declared model definition of Has Physical Capacitor Proton Parameters.
\end{proof}
\begin{definition}[Satisfies Uniform Parallel Plate Potential Law]
  \label{def:physics:phyx-mini-0866:phyxminiproblems-problemphyxmini0866-satisfiesuniformparallelplatepotentiallaw}
  \lean{PhyXMiniProblems.ProblemPhyXMini0866.SatisfiesUniformParallelPlatePotentialLaw}
  \uses{def:physics:phyx-mini-0866:phyxminiproblems-problemphyxmini0866-electricpotentialinvolts, def:physics:phyx-mini-0866:phyxminiproblems-problemphyxmini0866-capacitorprotonsetup}
  For ideal parallel plates, electric potential varies affinely across the gap. Thus the geometric midpoint has the average plate potential, and collision at a plate has that plate's potential. The turning-point inequalities encode that the shown reversal occurs strictly inside the gap
\end{definition}
\begin{proof}
  This is the declared model definition of Satisfies Uniform Parallel Plate Potential Law.
\end{proof}
\begin{definition}[Satisfies Electrostatic Energy Laws]
  \label{def:physics:phyx-mini-0866:phyxminiproblems-problemphyxmini0866-satisfieselectrostaticenergylaws}
  \lean{PhyXMiniProblems.ProblemPhyXMini0866.SatisfiesElectrostaticEnergyLaws}
  \uses{def:physics:phyx-mini-0866:phyxminiproblems-problemphyxmini0866-massinkilograms, def:physics:phyx-mini-0866:phyxminiproblems-problemphyxmini0866-chargemagnitudeincoulombs, def:physics:phyx-mini-0866:phyxminiproblems-problemphyxmini0866-speedinmeterspersecond, def:physics:phyx-mini-0866:phyxminiproblems-problemphyxmini0866-electricpotentialinvolts, def:physics:phyx-mini-0866:phyxminiproblems-problemphyxmini0866-energyinjoules, def:physics:phyx-mini-0866:phyxminiproblems-problemphyxmini0866-motionevent, def:physics:phyx-mini-0866:phyxminiproblems-problemphyxmini0866-capacitorprotonsetup}
  The nonrelativistic kinetic-energy law K = m v² / 2, electrostatic potential energy U = q V for a positive proton, and conservation of K + U throughout the ideal field. These are general laws over all three named events, not a specialization to the numerical answer
\end{definition}
\begin{proof}
  This is the declared model definition of Satisfies Electrostatic Energy Laws.
\end{proof}
\begin{definition}[Answer Choice]
  \label{def:physics:phyx-mini-0866:phyxminiproblems-problemphyxmini0866-answerchoice}
  \lean{PhyXMiniProblems.ProblemPhyXMini0866.AnswerChoice}
  Labels of the four displayed speed choices
\end{definition}
\begin{proof}
  This is the declared model definition of Answer Choice.
\end{proof}
\begin{definition}[displayed Speed In Meters Per Second]
  \label{def:physics:phyx-mini-0866:phyxminiproblems-problemphyxmini0866-displayedspeedinmeterspersecond}
  \lean{PhyXMiniProblems.ProblemPhyXMini0866.displayedSpeedInMetersPerSecond}
  \uses{def:physics:phyx-mini-0866:phyxminiproblems-problemphyxmini0866-answerchoice}
  Speed in metres per second printed beside an answer label
\end{definition}
\begin{proof}
  This is the declared model definition of displayed Speed In Meters Per Second.
\end{proof}
\begin{definition}[Rounds To Displayed Speed]
  \label{def:physics:phyx-mini-0866:phyxminiproblems-problemphyxmini0866-roundstodisplayedspeed}
  \lean{PhyXMiniProblems.ProblemPhyXMini0866.RoundsToDisplayedSpeed}
  \uses{def:physics:phyx-mini-0866:phyxminiproblems-problemphyxmini0866-speedinmeterspersecond, def:physics:phyx-mini-0866:phyxminiproblems-problemphyxmini0866-capacitorprotonsetup, def:physics:phyx-mini-0866:phyxminiproblems-problemphyxmini0866-answerchoice, def:physics:phyx-mini-0866:phyxminiproblems-problemphyxmini0866-displayedspeedinmeterspersecond}
  The printed values have three significant figures, so a speed agrees with a choice when it lies within half of the last displayed 1000 m/s increment
\end{definition}
\begin{proof}
  This is the declared model definition of Rounds To Displayed Speed.
\end{proof}
\begin{definition}[Is Closest Displayed Speed]
  \label{def:physics:phyx-mini-0866:phyxminiproblems-problemphyxmini0866-isclosestdisplayedspeed}
  \lean{PhyXMiniProblems.ProblemPhyXMini0866.IsClosestDisplayedSpeed}
  \uses{def:physics:phyx-mini-0866:phyxminiproblems-problemphyxmini0866-speedinmeterspersecond, def:physics:phyx-mini-0866:phyxminiproblems-problemphyxmini0866-capacitorprotonsetup, def:physics:phyx-mini-0866:phyxminiproblems-problemphyxmini0866-answerchoice, def:physics:phyx-mini-0866:phyxminiproblems-problemphyxmini0866-displayedspeedinmeterspersecond}
  A choice is at least as close to the physical collision speed as any other
\end{definition}
\begin{proof}
  This is the declared model definition of Is Closest Displayed Speed.
\end{proof}
\begin{lemma}[launch and collision potential readouts]
  \label{lem:physics:phyx-mini-0866:phyxminiproblems-problemphyxmini0866-launch-and-collision-potential-readouts}
  \lean{PhyXMiniProblems.ProblemPhyXMini0866.launch_and_collision_potential_readouts}
  \uses{def:physics:phyx-mini-0866:phyxminiproblems-problemphyxmini0866-electricpotentialinvolts, def:physics:phyx-mini-0866:phyxminiproblems-problemphyxmini0866-capacitorprotonsetup, def:physics:phyx-mini-0866:phyxminiproblems-problemphyxmini0866-matchescapacitorprotonscenario, def:physics:phyx-mini-0866:phyxminiproblems-problemphyxmini0866-matchessuppliedcapacitorfigureanddata, def:physics:phyx-mini-0866:phyxminiproblems-problemphyxmini0866-satisfiesuniformparallelplatepotentiallaw}
  The uniform-potential law places the launch midpoint at 250 V and the negative collision plate at 0 V. Both are derived relations
\end{lemma}
\begin{proof}
  The conclusion follows from the governing relations and figure readouts named in the statement of launch and collision potential readouts.
\end{proof}
\begin{lemma}[collision speed squared from energy conservation]
  \label{lem:physics:phyx-mini-0866:phyxminiproblems-problemphyxmini0866-collision-speed-squared-from-energy-conservation}
  \lean{PhyXMiniProblems.ProblemPhyXMini0866.collision_speed_squared_from_energy_conservation}
  \uses{def:physics:phyx-mini-0866:phyxminiproblems-problemphyxmini0866-massinkilograms, def:physics:phyx-mini-0866:phyxminiproblems-problemphyxmini0866-chargemagnitudeincoulombs, def:physics:phyx-mini-0866:phyxminiproblems-problemphyxmini0866-speedinmeterspersecond, def:physics:phyx-mini-0866:phyxminiproblems-problemphyxmini0866-electricpotentialinvolts, def:physics:phyx-mini-0866:phyxminiproblems-problemphyxmini0866-capacitorprotonsetup, def:physics:phyx-mini-0866:phyxminiproblems-problemphyxmini0866-hasphysicalcapacitorprotonparameters, def:physics:phyx-mini-0866:phyxminiproblems-problemphyxmini0866-satisfieselectrostaticenergylaws}
  Energy conservation determines the collision speed squared before any decimal rounding or comparison with the answer choices
\end{lemma}
\begin{proof}
  The conclusion follows from the governing relations and figure readouts named in the statement of collision speed squared from energy conservation.
\end{proof}
\begin{lemma}[collision speed rounds to choice D]
  \label{lem:physics:phyx-mini-0866:phyxminiproblems-problemphyxmini0866-collision-speed-rounds-to-choice-d}
  \lean{PhyXMiniProblems.ProblemPhyXMini0866.collision_speed_rounds_to_choice_D}
  \uses{def:physics:phyx-mini-0866:phyxminiproblems-problemphyxmini0866-capacitorprotonsetup, def:physics:phyx-mini-0866:phyxminiproblems-problemphyxmini0866-matchescapacitorprotonscenario, def:physics:phyx-mini-0866:phyxminiproblems-problemphyxmini0866-matchessuppliedcapacitorfigureanddata, def:physics:phyx-mini-0866:phyxminiproblems-problemphyxmini0866-hasstandardprotoncalibration, def:physics:phyx-mini-0866:phyxminiproblems-problemphyxmini0866-hasphysicalcapacitorprotonparameters, def:physics:phyx-mini-0866:phyxminiproblems-problemphyxmini0866-satisfiesuniformparallelplatepotentiallaw, def:physics:phyx-mini-0866:phyxminiproblems-problemphyxmini0866-satisfieselectrostaticenergylaws, def:physics:phyx-mini-0866:phyxminiproblems-problemphyxmini0866-roundstodisplayedspeed}
  With the standard proton calibration and the supplied readouts, the positive collision speed is within 500 m/s of 296000 m/s = 2.96 × 10⁵ m/s
\end{lemma}
\begin{proof}
  The conclusion follows from the governing relations and figure readouts named in the statement of collision speed rounds to choice D.
\end{proof}
% --- Archon declaration topology end ---
% --- Archon physics formalization source end ---
